% Pursuing Stacks (cotes 134-2 to 134-8): archive cote 134-6, site page 10 (PDF page 11), typescript page 329.
% typescript. Draft compared: ps5.tex:1832-1902.
% AI-transcribed by Claude Opus 5.5 (claude-opus-5-5, adaptive thinking, effort high), Message Batches API, batch ps3 (msgbatch_01XdFmuVVBAcgc7EFxdtCqsj), 2026-11-12; not yet checked by eye.
\scanpage{134-6/10}
\mspage{329}

where all functors in the diagram,  except for K and L,  are the canoni-%
cal ones familiar to us.  The question then is how to define a suitable
L,  such that the corresponding square (where  $K = L\circ Wh_{A}$ ) should commu-%
te up to (canonical ?) isomorphism.

There \uline{is} such an L,  whenever we have an\xout{test}  aspheric functor (10)
(where \hand{$\Delta$} may be replaced by one of it's twins),  using corresponding
\unsure{Semisimplicial|semisimplicial} chain complexes  -  but, as already remarked yesterday,
taking things this way is ``cheating''!  Conceivably too,  there are quite
general theorems asserting that a functor from a category $A^{\wedge}$ say to
a derived \xout{abelian}\inserted[typed]{additive} category such as Hotab  can be lifted to the category
of models (here \deleted{chain}\inserted[typed]{abelian} complexes)  it comes from,  and possibly even in a
way factoring through $A^{\wedge}_{ab}$ ?   I do'nt intend to dive into these questions
either,  but rather,  make a comment on a general method for constructing
certain functors
\hand{(21)} \[ L:\ A^{\wedge}_{ab} \longrightarrow Ch.(Ab) \quad , \]
(maybe not in a way to give rise to a commutative diagram (20)),  as
suggested by the standard chain complexes associated to the three types
of complexes,  using simpli\unsure{\hand{c}}es,  squares or hemispheres,  or multicom-%
plexes (using products of the standard test categories).  Writing
\[ A^{\wedge}_{ab} \hand{\xrightarrow{\sim}} \underline{Hom}(A^{\unsure{\hand{\circ}}}, (Ab)) \quad , \]
we remark that the standard constructions of chain complexes associated
to semisimmlicial (say) complexes of abelian groups,  makes a sense not
only for such abelian complexes,  but more generally for \uline{complexes with
values in} any \uline{addit\xout{\unsure{t}}ive category},  M say.   This induces us to look
more generally,  for any such M,  for a functor
\hand{(2\unsure{2})} \[ L_{M} :\ \underline{Hom}(A^{O}, M) \longrightarrow Ch.(M) \quad , \]
in a way \uline{compatible with additive functors}
\[ M \to M' \]
(in the obvious sense of the word).\hand{//} Now for any category B (here A\xout{\illegible{1}}\hand{$^{O}$} )
one can define an "\uline{envelopping additive category}"  Add(B),  together
with a canonical functor
\hand{(2\unsure{3})} \[ B \longrightarrow Add(B) \quad , \]
which \xout{\unsure{should}} is "2-universal"  for all possible functors of B into any
additive category M.  More specifically,  for any such M,  the corres-%

\margin{10}
